%! AP Statistics — Unit 4, Lesson 4.4 — Group Activity
\documentclass[10pt]{article}
\usepackage{apstats-article}
\usepackage{apstats-boxes}

\begin{document}

\pageheader{Unit 4, Lesson 4.4}{Group Activity: Are These Events Mutually Exclusive?}
\namepartnerperiod

\vspace{0.1in}

\begin{scenariobox}[Context]{sky}
\small
A large high school surveyed 200 students about their after-school activities and
class standing. The two-way table below shows the results.

\vspace{6pt}
\renewcommand{\arraystretch}{1.4}
\begin{tabularx}{\linewidth}{@{} l *{4}{X} r @{}}
 & \textbf{Freshman} & \textbf{Sophomore} & \textbf{Junior} & \textbf{Senior} & \textbf{Total} \\
\hline
\textbf{Sports}    & 18 & 22 & 20 & 15 & 75  \\
\textbf{Arts/Music}& 12 & 10 & 14 & 9  & 45  \\
\textbf{Neither}   & 20 & 18 & 22 & 20 & 80  \\
\hline
\textbf{Total}     & 50 & 50 & 56 & 44 & 200 \\
\end{tabularx}
\end{scenariobox}

\vspace{0.12in}

% ── TIER R ────────────────────────────────────────────────────────────────────
\begin{tcolorbox}[colback=greenbg, colframe=greenacc, boxrule=0.9pt, arc=2mm,
    left=3mm, right=3mm, top=2mm, bottom=2mm,
    title={\bfseries\small Tier R --- Remediate: Identifying Mutually Exclusive Events},
    fonttitle=\small]
\small
For each pair of events, circle \textbf{ME} (mutually exclusive) or \textbf{Not ME}.
A brief reason is given to help you.

\vspace{6pt}
\begin{enumerate}[leftmargin=1.5em, itemsep=8pt]
  \item $A = \{\text{student is a Freshman}\}$, $B = \{\text{student is a Senior}\}$.\\
        Can a student be both a freshman and a senior at the same time?
        \blank{3cm}\\
        \textbf{ME} \quad / \quad \textbf{Not ME}

  \item $A = \{\text{student is in Sports}\}$, $B = \{\text{student is in Arts/Music}\}$.\\
        Can a student be in both activities? (Look at the table --- are there rows for both?)
        \blank{3cm}\\
        \textbf{ME} \quad / \quad \textbf{Not ME}

  \item $A = \{\text{student is in Sports}\}$, $B = \{\text{student is a Junior}\}$.\\
        Can a student play sports and be a junior? \blank{3cm}\\
        \textbf{ME} \quad / \quad \textbf{Not ME}

  \item $A = \{\text{student is in Arts/Music}\}$, $B = \{\text{student is in Neither}\}$.\\
        \textbf{ME} \quad / \quad \textbf{Not ME}
\end{enumerate}
\end{tcolorbox}

\vspace{0.1in}

% ── TIER A ────────────────────────────────────────────────────────────────────
\begin{tcolorbox}[colback=sky, colframe=navy, boxrule=0.9pt, arc=2mm,
    left=3mm, right=3mm, top=2mm, bottom=2mm,
    title={\bfseries\small Tier A --- Approaching Proficiency: Classify, Justify, and Compute},
    fonttitle=\small]
\small
Use the table above to answer each question. Show your reasoning.

\vspace{6pt}
\begin{enumerate}[leftmargin=1.5em, itemsep=10pt]
  \item Let $A = \{\text{Sophomore}\}$ and $B = \{\text{Sports}\}$. One student is chosen
        at random.
        \begin{enumerate}[label=(\alph*), itemsep=5pt]
          \item Are $A$ and $B$ mutually exclusive? Justify in context.\\
                \writeline\writeline
          \item Find $P(A \cap B)$. Show your setup. \blank{5cm}
        \end{enumerate}

  \item Let $A = \{\text{Freshman}\}$ and $B = \{\text{Senior}\}$.
        \begin{enumerate}[label=(\alph*), itemsep=5pt]
          \item Are $A$ and $B$ mutually exclusive? Justify in context.\\
                \writeline\writeline
          \item Find $P(A \cap B)$. \blank{4cm}
        \end{enumerate}

  \item Let $A = \{\text{Sports}\}$ and $B = \{\text{Arts/Music}\}$.
        Based on the table structure (each row is a separate activity), are students
        assumed to be in at most one activity?
        \begin{enumerate}[label=(\alph*), itemsep=5pt]
          \item Are $A$ and $B$ mutually exclusive? Justify.\\
                \writeline\writeline
          \item Find $P(A \cap B)$ and $P(A)$. \blank{5cm}
        \end{enumerate}
\end{enumerate}
\end{tcolorbox}

\vspace{0.1in}

% ── TIER E ────────────────────────────────────────────────────────────────────
\begin{tcolorbox}[colback=goldbg, colframe=goldacc, boxrule=0.9pt, arc=2mm,
    left=3mm, right=3mm, top=2mm, bottom=2mm,
    title={\bfseries\small Tier E --- Extension: Creating and Critiquing},
    fonttitle=\small]
\small
\begin{enumerate}[leftmargin=1.5em, itemsep=10pt]
  \item \textbf{Create your own:} Using the same table, define two events that \emph{are}
        mutually exclusive and two events that are \emph{not} mutually exclusive.
        Write a complete sentence justification for each.

        \vspace{4pt}
        ME pair: $A = $ \blank{5cm}, $B = $ \blank{5cm}\\
        Justification: \writeline\writeline

        \vspace{4pt}
        Not ME pair: $A = $ \blank{5cm}, $B = $ \blank{5cm}\\
        Justification: \writeline\writeline

  \item \textbf{Critique this claim:} A classmate says: ``Events $A$ and $B$ are
        mutually exclusive because $P(A)$ and $P(B)$ are both small.''
        Is this reasoning correct? Explain the error.\\
        \writeline\writeline\writeline

  \item \textbf{Probability challenge:} A student says: ``If two events are mutually exclusive,
        then neither can have a high probability.'' Provide a counterexample that shows this
        is false, then explain.\\
        \writeline\writeline\writeline
\end{enumerate}
\end{tcolorbox}

\end{document}
